\documentclass[11pt,a4paper]{article}
\usepackage[utf8]{inputenc}
\usepackage[vietnamese]{babel}
\usepackage{amsmath,amssymb}
\usepackage{geometry}
\geometry{a4paper, total={180mm,260mm}, left=15mm, right=15mm, top=18mm, bottom=18mm}

\usepackage{tikz}
\usetikzlibrary{shapes.geometric, arrows.meta, positioning, calc}
\tikzset{arr/.style={-{Stealth[scale=0.8]}}}
\usepackage{tcolorbox}
\tcbuselibrary{skins,breakable}
\usepackage{fontawesome}
\usepackage{fancyhdr}
\usepackage{enumitem}

\pagestyle{fancy}
\fancyhf{}
\renewcommand{\headrulewidth}{0.5pt}
\renewcommand{\footrulewidth}{0.5pt}
\fancyhead[L]{\small\bfseries\color{violet!85!black}\faBook\ SỔ TAY KIẾN THỨC TOÁN 2}
\fancyhead[R]{\small\bfseries\color{violet!85!black}Lớp Toán Thầy Lĩnh -- loptoanthaylinh.vn}
\fancyfoot[C]{\small\color{gray}-- Sổ tay kiến thức: Bài 09. Bài toán về thêm, bớt một số đơn vị --}

\setlength{\parindent}{0pt}
\setlength{\parskip}{0pt}

\begin{document}

% --- BANNER BÀI HỌC ---
\begin{tcolorbox}[
  enhanced,
  colback=violet!8!white,
  colframe=violet!70!black,
  arc=3mm, boxrule=1.2pt,
  left=3mm, right=3mm, top=2mm, bottom=2mm,
  halign=center
]
  {\LARGE\bfseries\color{violet!90!black}\faBookmark\ BÀI 09: BÀI TOÁN VỀ THÊM, BỚT MỘT SỐ ĐƠN VỊ}\\[1mm]
  {\itshape\footnotesize\color{black!75} Quy tắc nhận diện từ khóa và phương pháp giải bài toán thêm vào, bớt đi trên 1 trang A4}
\end{tcolorbox}

\vspace{-2mm}

\noindent
\begin{minipage}[t]{0.49\textwidth}\vspace{0pt}
\begin{tcolorbox}[
  enhanced,
  colback=violet!4!white,
  colframe=violet!75!black,
  arc=2mm, boxrule=0.8pt,
  left=2mm, right=2mm, top=1.8mm, bottom=1.8mm,
  title={\footnotesize\bfseries\color{white}\faKey\ 1. Bài toán về "thêm"}
]
{\footnotesize
\begin{itemize}\setlength{\itemsep}{1pt}\setlength{\leftskip}{-2.5mm}
  \item \textbf{Ý nghĩa:} "Thêm" nghĩa là gộp vào, nhận thêm, mua thêm, bay đến,... làm cho số lượng tăng lên.
  \item \textbf{Phép tính giải:}\\
  Ta thực hiện \textbf{phép tính cộng} ($+$):
  $$\mathbf{\text{Lúc sau} = \text{Lúc đầu} + \text{Phần thêm}}$$
  \item \textbf{Ví dụ:} Có $8$ quả táo, mẹ cho thêm $5$ quả. Có tất cả: $8 + 5 = 13$ quả táo.
\end{itemize}
}
\end{tcolorbox}
\end{minipage}
\hfill
\begin{minipage}[t]{0.49\textwidth}\vspace{0pt}
\begin{tcolorbox}[
  enhanced,
  colback=teal!4!white,
  colframe=teal!75!black,
  arc=2mm, boxrule=0.8pt,
  left=2mm, right=2mm, top=1.8mm, bottom=1.8mm,
  title={\footnotesize\bfseries\color{white}\faSliders\ 2. Bài toán về "bớt"}
]
{\footnotesize
\begin{itemize}\setlength{\itemsep}{1pt}\setlength{\leftskip}{-2.5mm}
  \item \textbf{Ý nghĩa:} "Bớt" nghĩa là lấy đi, cho đi, bán đi, ăn mất, bay đi,... làm cho số lượng giảm đi.
  \item \textbf{Phép tính giải:}\\
  Ta thực hiện \textbf{phép tính trừ} ($-$):
  $$\mathbf{\text{Còn lại} = \text{Lúc đầu} - \text{Phần bớt}}$$
  \item \textbf{Ví dụ:} Có $15$ con chim, bay đi $6$ con. Còn lại: $15 - 6 = 9$ con chim.
\end{itemize}
}
\end{tcolorbox}
\end{minipage}

\vspace{-1.5mm}
\begin{tcolorbox}[
  enhanced,
  colback=orange!4!white,
  colframe=orange!80!black,
  arc=2mm, boxrule=0.8pt,
  left=2mm, right=2mm, top=1.5mm, bottom=1.5mm,
  title={\footnotesize\bfseries\color{white}\faSitemap\ 3. Sơ đồ phân nhánh từ khóa nhận diện phép tính}
]
\centering
\begin{tikzpicture}[
  node distance=0.8cm,
  box/.style={rectangle, rounded corners=2pt, draw=violet!80!black, fill=violet!15, text width=2.4cm, align=center, font=\scriptsize\bfseries, inner sep=3pt},
  branch/.style={rectangle, rounded corners=2pt, draw=#1!80!black, fill=#1!10, text width=3.4cm, align=center, font=\scriptsize, inner sep=2.5pt},
  arr/.style={-{Stealth[scale=0.8]}, thick, #1!80!black}
]
  \node[box] (root) {BÀI TOÁN\\THÊM -- BỚT};
  \node[branch=teal, left=0.6cm of root] (b1) {\textbf{BÀI TOÁN THÊM}\\Từ khóa: Thêm, mua, gộp\\$\implies$ \textbf{Làm phép CỘNG (+)}};
  \node[branch=red, right=0.6cm of root] (b2) {\textbf{BÀI TOÁN BỚT}\\Từ khóa: Bớt, cho, bay đi\\$\implies$ \textbf{Làm phép TRỪ (-)}};

  \draw[arr=teal] (root) -- (b1);
  \draw[arr=red] (root) -- (b2);
\end{tikzpicture}
\end{tcolorbox}

\vspace{-1.5mm}
\noindent
\begin{minipage}[t]{0.49\textwidth}\vspace{0pt}
\begin{tcolorbox}[
  enhanced,
  colback=green!4!white,
  colframe=green!60!black,
  arc=2mm, boxrule=0.8pt,
  left=2mm, right=2mm, top=1.8mm, bottom=1.8mm,
  title={\footnotesize\bfseries\color{white}\faCheckCircle\ 4. Ví dụ thực hành chuẩn mực}
]
{\footnotesize
\textbf{Bài toán mẫu:}\\
Trên cành cây có $9$ con chim, một lát sau có thêm $4$ con chim bay đến đậu cùng. Hỏi trên cành cây có tất cả bao nhiêu con chim?\\
\textit{Bài giải:}\\
Trên cành cây có tất cả số con chim là:
$$9 + 4 = 13\text{ (con chim)}$$
Đáp số: $13$ con chim.
}
\end{tcolorbox}
\end{minipage}
\hfill
\begin{minipage}[t]{0.49\textwidth}\vspace{0pt}
\begin{tcolorbox}[
  enhanced,
  colback=red!4!white,
  colframe=red!75!black,
  arc=2mm, boxrule=0.8pt,
  left=2mm, right=2mm, top=1.8mm, bottom=1.8mm,
  title={\footnotesize\bfseries\color{white}\faExclamationTriangle\ 5. Mẹo tránh bẫy trắc nghiệm}
]
{\footnotesize
\begin{itemize}\setlength{\itemsep}{1pt}\setlength{\leftskip}{-2.5mm}
  \item \textbf{Bẫy đọc không kĩ câu hỏi:}\\
  Đề bài cho thêm vào nhưng lại hỏi "còn lại", hoặc đề cho bớt đi lại hỏi "tất cả". Phải đọc kỹ từ khóa!
  \item \textbf{Phân biệt "thêm vào" vs "nhiều hơn":}\\
  - Thêm vào: tác động làm tăng số lượng của cùng 1 đối tượng.\\
  - Nhiều hơn: so sánh chênh lệch giữa 2 đối tượng khác nhau.
\end{itemize}
}
\end{tcolorbox}
\end{minipage}


\end{document}
